\chapter{Deret Pangkat}\label{ch:b2:powerseries}

Deret pangkat adalah deret fungsi yang paling baik kelakuannya dalam
matematika: di dalam cakram kekonvergenannya ia konvergen
normal pada himpunan \omterm{def:b2:metric:compact}{kompak}, boleh diturunkan dan diintegralkan suku
demi suku tanpa berpikir dua kali, dan jumlahnya --- yakni \emph{\omterm{def:b2:powerseries:analytic}{fungsi
analitik}} --- ditentukan oleh koefisiennya. Bab ini
membuktikan seluruh paket itu lalu merebut kembali, secara jujur, setiap
deret Taylor pada jilid Tahun ke-1; sedangkan \omterm{ex:b2:powerseries:fibonacci}{fungsi pembangkit}
menutupnya dengan panen aljabar.

\section{Jari-jari kekonvergenan}

\begin{lemma}[Abel]\label{lem:b2:powerseries:abel}
Jika barisan $(a_n z_0^n)$ terbatas untuk suatu $z_0 \neq 0$,
maka $\sum a_n z^n$ konvergen \omterm{def:b2:series:def}{mutlak} untuk setiap $\abs z < \abs
{z_0}$, dan konvergen normal pada tiap cakram $\abs z \leq r < \abs{z_0}$.
\end{lemma}

\begin{proof}
Dengan $\abs{a_n z_0^n} \leq M$ dan $\abs z \leq r$:
\[
\abs{a_n z^n} = \abs{a_n z_0^n}\,\Bigl|\frac{z}{z_0}\Bigr|^n
\leq M\Bigl(\frac{r}{\abs{z_0}}\Bigr)^{\!n},
\]
yakni batas geometri yang konvergen, \omterm{def:b2:funcseq:def}{seragam} pada cakramnya.
\end{proof}

\begin{definition}[Jari-jari kekonvergenan]\label{def:b2:powerseries:radius}
\emph{Jari-jari kekonvergenan}\index{jari-jari kekonvergenan} bagi
$\sum a_n z^n$ adalah
\[
R = \sup\{r \geq 0 : (a_n r^n) \text{ terbatas}\} \in
\intcc{0}{+\infty} .
\]
Menurut \cref{lem:b2:powerseries:abel}: kekonvergenan \omterm{def:b2:series:def}{mutlak} untuk
$\abs z < R$ (dan normal pada subcakram \omterm{def:b2:metric:compact}{kompak}), sedangkan kedivergenan
--- bahkan suku yang tak terbatas --- untuk $\abs z > R$. Pada lingkaran
batasnya, apa saja bisa terjadi (\cref{exo:b2:powerseries:2}). Dalam
praktik $R$ dihitung lewat uji rasio d'Alembert pada $\abs{a_n}\abs z^n$
atau lewat pembandingan.
\end{definition}

\begin{example}[Sebuah jari-jari tanpa uji rasio]\label{ex:b2:powerseries:sinn}
Berapa jari-jari $\sum \sin(n)\,z^n$? Rasio
$\abs{\sin(n+1)/\sin n}$ tak berlimit, tetapi definisinya bekerja
secara langsung. \emph{$R \geq 1$:} sebab $\abs{\sin n} \leq 1$, jadi
$(\sin n\cdot r^n)$ terbatas untuk setiap $r < 1$ --- bahkan untuk $r =
1$. \emph{$R \leq 1$:} cukuplah bahwa $\sin n \not\to 0$.
Andaikan $\sin n \to 0$; maka rumus penjumlahan
\[
\sin(n+1) = \sin n\cos 1 + \cos n\sin 1
\]
akan memaksa $\cos n \to 0$ (selesaikan untuk $\cos n$, sebab $\sin 1
\neq 0$), yang bertentangan dengan $\sin^2 n + \cos^2 n = 1$. Jadi suku
$\sin(n)\,1^n$ tidak menuju $0$: sehingga deretnya divergen di $z =
1$, dan $R \leq 1$. Kesimpulannya: $R = 1$. Pelajaran penutupnya:
jari-jarinya adalah pernyataan tentang keterbatasan $\abs{a_n}r^n$ ---
jadi limit rasio tak pernah diperlukan, dan hujah keterbatasan
menyelesaikan kasus yang tak tersentuh uji rasio (bandingkan dengan
koefisien berayun pada \cref{exo:b2:powerseries:1}).
\end{example}

\begin{proposition}[Operasi]\label{prop:b2:powerseries:operations}
Misalkan $\sum a_nz^n$ dan $\sum b_nz^n$ berjari-jari $R_a, R_b$. Maka,
untuk $\abs z < \min(R_a, R_b)$:
\[
\sum (a_n + b_n)z^n = \sum a_nz^n + \sum b_nz^n,
\qquad
\Bigl(\sum a_nz^n\Bigr)\Bigl(\sum b_nz^n\Bigr) = \sum c_n z^n,
\quad c_n = \sum_{k=0}^{n} a_kb_{n-k},
\]
dengan kedua deretnya berjari-jari $\geq \min(R_a, R_b)$. (Hasil kalinya
adalah hasil kali Cauchy, yang sah berkat kekonvergenan \omterm{def:b2:series:def}{mutlak} dan
\cref{thm:b2:series:fubini}.)
\end{proposition}

\begin{proof}
Rumus jumlahnya tak lain kelinearan deret konvergen, dan
$(a_n + b_n)r^n$ terbatas setiap kali $a_nr^n$ dan $b_nr^n$
terbatas: jadi jari-jarinya $\geq \min(R_a, R_b)$. Untuk hasil kalinya,
tetapkan $\abs z < \min(R_a, R_b)$: kedua deretnya konvergen
\emph{\omterm{def:b2:series:def}{mutlak}} di sana (\cref{lem:b2:powerseries:abel}), jadi keluarga
berindeks ganda $(a_kz^k\,b_lz^l)_{k,l}$ \omterm{def:b2:series:summable}{terjumlahkan}, sehingga
\cref{thm:b2:series:fubini} mengizinkan pengelompokan apa pun. Setelah
dikelompokkan menurut $k + l = n$:
\[
\Bigl(\sum_k a_kz^k\Bigr)\Bigl(\sum_l b_lz^l\Bigr)
= \sum_{n\geq0}\Bigl(\sum_{k+l=n}a_kb_l\Bigr)z^n
= \sum_{n\geq0}c_nz^n ,
\]
yang konvergen \omterm{def:b2:series:def}{mutlak} untuk setiap $z$ semacam itu: jadi deret hasil
kalinya juga berjari-jari $\geq \min(R_a, R_b)$.
\end{proof}

\begin{example}[Sebuah kuadrat Cauchy, diperiksa silang]\label{ex:b2:powerseries:cauchysquare}
Kuadratkan deret geometrinya: untuk $\abs x < 1$, koefisien
$x^n$ pada $\bigl(\sum x^k\bigr)^2$ adalah $c_n = \sum_{k+l=n}
1\cdot1 = n + 1$, jadi
\[
\frac{1}{(1-x)^2} = \sum_{n\geq0}(n+1)\,x^n .
\]
Periksa silang lewat penurunan suku demi suku
(\cref{thm:b2:powerseries:calculus} di bawah): menurunkan
$\frac{1}{1-x} = \sum x^n$ memberi $\frac{1}{(1-x)^2} = \sum
nx^{n-1} = \sum(n+1)x^n$ --- yakni deret yang sama lewat dua
mekanisme yang tak berkaitan. Pelajaran penutupnya: ketika sebuah
kesamaan koefisien tampak misterius, salah satu dari kedua mesin ini
(konvolusi atau penurunan) biasanya menghasilkannya dalam satu baris;
sedangkan pertanyaan soal akhir pekan tentang $\sum\binom{2k}k\binom{2n-2k}{n-k} = 4^n$
menjalankan mesin konvolusinya dengan tenaga penuh.
\end{example}

\begin{example}[Mengalikan dengan $\frac{1}{1-x}$ menjumlahkan koefisiennya]\label{ex:b2:powerseries:partialsums}
\omterm{thm:b2:series:fubini}{Hasil kali Cauchy} terhadap deret geometri punya makna yang mudah
diingat: untuk sembarang $\sum a_nx^n$ berjari-jari $R > 0$ dan $\abs x <
\min(R, 1)$,
\[
\frac{1}{1-x}\sum_{n\geq0}a_nx^n
= \sum_{n\geq0}\Bigl(\sum_{k=0}^{n}a_k\Bigr)x^n :
\]
jadi mengalikan dengan $\frac{1}{1-x}$ mengganti koefisiennya dengan
jumlah parsialnya (yakni mengonvolusikannya dengan barisan yang semua
sukunya satu). Contohnya: $\dfrac{\eu^x}{1-x} = \sum_n s_n x^n$ dengan
$s_n = \sum_{k\leq n}\frac{1}{k!}$, yakni jumlah parsial $\eu$ ---
bandingkan dengan \cref{exo:b2:powerseries:11}, tempat hasil kali yang
sama dengan $\eu^{-x}$ menyandikan cacah permutasi tanpa titik tetap.
Pelajaran penutupnya: operasi pada deret pangkat adalah operasi pada
barisan koefisien yang menyamar (kalikan dengan $\frac1{1-x}$:
menjumlahkan; kalikan dengan $x$: menggeser; turunkan: mengalikan dengan
$n$ lalu menggeser) --- yakni kamus yang kelak dibaca fasih oleh bab
\omterm{ex:b2:powerseries:fibonacci}{fungsi pembangkit}.
\end{example}

\section{Keteraturan jumlahnya}

\begin{theorem}[Kalkulus suku demi suku]\label{thm:b2:powerseries:calculus}
Misalkan $f(x) = \sum_{n\geq0} a_n x^n$ berjari-jari $R > 0$ (dengan
peubah real $x \in \intoo{-R}{R}$).
\begin{enumerate}
  \item Deret turunannya $\sum n\,a_n x^{n-1}$ punya jari-jari $R$
        yang \emph{sama}, dan $f$ berkelas $C^1$ dengan $f'(x) =
        \sum_{n \geq 1} n a_n x^{n-1}$. Setelah diiterasikan, $f$
        berkelas $C^\infty$ dan
        \[
        a_n = \frac{f^{(n)}(0)}{n!} :
        \]
        jadi koefisien sebuah deret pangkat bersifat tunggal (dua deret
        yang jumlahnya sama di dekat $0$ punya koefisien yang sama).
  \item Antiturunan suku demi suku: $\sum \frac{a_n}{n+1}x^{n+1}$
        berjari-jari $R$ dengan turunan $f$.
\end{enumerate}
\end{theorem}

\begin{proof}
\emph{Jari-jarinya sama:} jika $(a_nr^n)$ terbatas dan $r' < r$, maka
$n\abs{a_n} r'^{\,n-1} = \frac{n}{r'}\abs{a_nr^n}
\bigl(\frac{r'}{r}\bigr)^n$ terbatas (bahkan $\to 0$: sebab geometri
mengalahkan $n$), jadi $R' \geq R$; sebaliknya $\abs{a_n x^n} \leq \abs x
\cdot n\abs{a_n}\abs x^{n-1}$ memberi $R \geq R'$.

\emph{Penurunannya:} pada $\intcc{-r}{r}$ dengan $r < R$, deret
turunannya konvergen normal (sebab $n\abs{a_n}r^{n-1}$ \omterm{def:b2:series:summable}{terjumlahkan}
menurut perhitungan jari-jarinya); sedangkan deret aslinya konvergen di
$x = 0$: jadi teorema penurunan untuk deret
(\cref{thm:b2:funcseq:seriestransfer}) berlaku pada tiap ruas semacam
itu, sehingga pada $\intoo{-R}{R}$. Setelah diiterasikan $k$ kali lalu
dinilai di $0$: secara gamblang, deret turunan ke-$k$-nya adalah
\[
f^{(k)}(x) = \sum_{n\geq k} n(n-1)\cdots(n-k+1)\,a_n\,x^{n-k},
\]
dan di $x = 0$ setiap suku dengan $n > k$ lenyap, sehingga tersisa
hanya suku tetapnya $k(k-1)\cdots1\cdot a_k$: jadi $f^{(k)}(0) =
k!\,a_k$. Ketunggalan koefisiennya menyusul: dua deret pangkat
yang jumlahnya sama di dekat $0$ punya turunan yang sama di $0$,
sehingga punya $a_k$ yang sama. Untuk antiturunannya: jari-jarinya sama
lewat perhitungan yang sama, lalu turunkan suku demi suku kembali.
\end{proof}

\begin{example}[Menilai sebuah deret di sebuah titik]\label{ex:b2:powerseries:evaluate}
Berapa $\sum_{n\geq1}\dfrac{n^2}{2^n}$? Ia tak lain jumlah
$\sum n^2x^n$ pada \cref{exo:b2:powerseries:3} yang dinilai
\emph{di dalam} cakramnya, di $x = \frac12 < 1 = R$, tempat setiap
manipulasi yang dipakai untuk menurunkan bentuk tertutupnya memang sah:
\[
\sum_{n\geq1} n^2x^n = \frac{x(1+x)}{(1-x)^3}
\quad\Longrightarrow\quad
\sum_{n\geq1}\frac{n^2}{2^n}
= \frac{\frac12\cdot\frac32}{(\frac12)^3}
= \frac{3/4}{1/8} = 6 .
\]
Mesin yang sama, tombol yang lain: $x = \frac13$ memberi
$\sum\frac{n^2}{3^n} = \frac{\frac13\cdot\frac43}{(2/3)^3} =
\frac32$. Pelajaran penutupnya: sebuah kesamaan deret pangkat adalah
mesin, bukan satu rumus tunggal --- sebab satu penurunan menghargai
setiap deret numerik $\sum n^2q^n$ sekaligus, untuk setiap $\abs q < 1$;
dan begitulah bab \omterm{ex:b2:powerseries:fibonacci}{fungsi pembangkit} kelak menghitung
nilai harapan dan ragam secara borongan.
\end{example}

\begin{example}[Yang klasik, kali ini secara jujur]\label{ex:b2:powerseries:classics}
$\displaystyle\frac{1}{1 - x} = \sum x^n$ ($R = 1$); lalu setelah
diintegralkan suku demi suku (\cref{thm:b2:powerseries:calculus}~(2)):
\[
-\ln(1 - x) = \sum_{n\geq1} \frac{x^n}{n},
\qquad
\arctan x = \sum_{n \geq 0} \frac{(-1)^n x^{2n+1}}{2n+1}
\quad (\abs x < 1),
\]
dengan yang kedua lewat dua langkah: sulihkan $-x^2$ pada deret
geometrinya untuk memperoleh $\frac{1}{1+x^2} = \sum(-1)^nx^{2n}$
(berjari-jari $1$, sebab $\abs{x^2} < 1 \iff \abs x < 1$), lalu ambil
antiturunan suku demi suku yang nol di $0$; sebab kedua ruasnya
antiturunan fungsi yang sama dengan nilai yang sama di $0$,
sehingga keduanya sama pada $\intoo{-1}{1}$. Dan untuk $\exp$:
deret $E(x) = \sum \frac{x^n}{n!}$ ($R = \infty$) memenuhi
$E' = E$ dan $E(0) = 1$ lewat penurunan suku demi suku, jadi $E = \exp$
berkat ketunggalan Tahun ke-1. Karenanya setiap ``uraian baku'' pada
jilid Tahun ke-1 kini menjadi teorema tentang deret pangkat penuhnya.
\end{example}

\begin{example}[Sebuah logaritma dihitung dari dalam cakramnya]\label{ex:b2:powerseries:lntwo}
Menilai $-\ln(1-x) = \sum\frac{x^n}{n}$ di titik dalam
$x = \frac12$ memberi
\[
\sum_{n\geq1}\frac{1}{n\,2^n} = \ln 2 ,
\]
yakni penyajian $\ln 2$ yang konvergen cepat (sepuluh suku saja sudah
memberi $0.69306\ldots$ terhadap $\ln 2 = 0.69314\ldots$), jauh
lebih baik daripada deret berselang-seling $1 - \frac12 + \frac13 -
\dots$ yang hanya tersedia di perbatasannya. Pelajaran penutupnya:
setiap kali sebuah konstanta terjangkau baik di tepi maupun tegas
di dalam cakramnya, bagian dalamnya menang secara numerik --- yakni
peluruhan geometri melawan peluruhan harmonik.
\end{example}

\begin{example}[Penurunan memelihara jari-jarinya, bukan perbatasannya]\label{ex:b2:powerseries:boundaryloss}
Deret $\sum_{n\geq1}\frac{x^n}{n^2}$ berjari-jari $1$ dan
konvergen di \emph{kedua} ujungnya (yakni $\sum\frac1{n^2}$ dan
kembaran berselang-selingnya). Deret turunannya,
\[
\sum_{n\geq1}\frac{x^{n-1}}{n} ,
\]
punya jari-jari $1$ yang sama --- sebagaimana dijamin
\cref{thm:b2:powerseries:calculus} --- tetapi kini
divergen di $x = 1$ (yakni deret harmonik) sedangkan masih konvergen
di $x = -1$ (yakni yang berselang-seling). Satu penurunan lagi
menghasilkan $\sum_{n\geq2}\frac{n-1}{n}x^{n-2}$, yang divergen di
kedua ujungnya (sebab sukunya tidak menuju $0$). Pelajaran penutupnya:
tiap penurunan mengalikan koefisiennya dengan $n$, dan hal itu
tak pernah menggeser jari-jarinya (sebab geometri mengalahkan
polinomial) tetapi memakan satu orde peluruhan di perbatasannya; jadi
kalkulus suku demi suku adalah olahraga dalam ruangan, dan apa pun yang
terjadi di tepinya wajib diperiksa ulang --- dan teori Abel--Tauber pada
soal akhir pekan tak lain pemeriksaan ulang itu.
\end{example}

\begin{example}[Memecah sebuah deret menurut sisanya --- dikerjakan tuntas]\label{ex:b2:powerseries:quarter}
Hitunglah $f(x) = \sum_{n\geq0} \dfrac{x^{4n}}{(4n)!}$ dalam bentuk
tertutup. Baik $\cosh x = \sum \frac{x^{2m}}{(2m)!}$ maupun $\cos x =
\sum \frac{(-1)^m x^{2m}}{(2m)!}$ berjari-jari $\infty$, jadi
rata-ratanya boleh dihitung suku demi suku:
\[
\frac{\cosh x + \cos x}{2}
= \sum_{m\geq0}\frac{1 + (-1)^m}{2}\,\frac{x^{2m}}{(2m)!}
= \sum_{m \text{ genap}}\frac{x^{2m}}{(2m)!}
= \sum_{n\geq0}\frac{x^{4n}}{(4n)!} = f(x) .
\]
Tapis $\frac{1+(-1)^m}{2}$ menyimpan tepat $m$ yang genap: dan inilah
awatara real bagi tapis akar-satuan (versi kompleksnya, dengan $\iu^n$,
menyarikan sisa modulo $4$ dalam sekali tarikan). Pemeriksaan
penutupnya: $f$ menyelesaikan $f'''' = f$ dengan $f(0) = 1$ dan
$f'(0) = f''(0) = f'''(0) = 0$ --- turunkan deretnya empat
kali (\cref{thm:b2:powerseries:calculus}) lalu perhatikan ia
memproduksi dirinya sendiri; sedangkan $\frac{\cosh + \cos}{2}$
memenuhi data yang sama.
\end{example}

\begin{definition}[Fungsi analitik]\label{def:b2:powerseries:analytic}
Fungsi $f$ disebut \emph{analitik}\index{fungsi analitik} di $x_0$
apabila ia jumlah sebuah deret pangkat dalam $(x - x_0)$ pada sebuah
persekitaran; dan analitik pada sebuah interval apabila demikian di
setiap titiknya. Jumlah deret pangkat bersifat analitik
di dalam cakramnya (lewat penataan ulang uraiannya --- yang diterima
tanpa bukti pada tingkat ini untuk pemusatan ulang, sebab kasus $x_0 =
0$ tak lain \cref{thm:b2:powerseries:calculus}). Analitik mengakibatkan
$C^\infty$; sedangkan konversnya \emph{gagal}: lihat fungsi datar
$\eu^{-1/x^2}$ (\cref{exo:b2:powerseries:7}).
\end{definition}

\begin{example}[Pemusatan ulang, dan jari-jari sebagai jarak]\label{ex:b2:powerseries:recenter}
Uraikan $f(x) = \frac{1}{1-x}$ di sekitar $x_0 = \frac12$: dengan
menulis $x = \frac12 + h$,
\[
\frac{1}{1 - x} = \frac{1}{\frac12 - h}
= \frac{2}{1 - 2h}
= \sum_{n\geq0} 2^{n+1}\,h^n
= \sum_{n\geq0} 2^{n+1}\Bigl(x - \frac12\Bigr)^{\!n},
\]
yang sahih untuk $\abs{2h} < 1$, yakni $\abs{x - \frac12} < \frac12$.
Jari-jari barunya persis sama dengan jarak dari pusat barunya ke
kesingularan $x = 1$: jadi pemusatan ulang menciutkan (atau membesarkan)
cakramnya agar pas dengan halangan terdekatnya. Pelajaran penutupnya:
inilah gambaran di balik definisi keanalitikan --- satu
fungsi, banyak deret pangkat lokal, yang masing-masing hidup pada cakram
terbesar yang menghindari masalahnya; dan jilid Tahun ke-3 mengubah
heuristik ``jari-jari $=$ jarak ke kesingularan kompleks terdekat''
menjadi sebuah teorema.
\end{example}

\begin{remark}[Jebakan yang sering muncul]
\emph{(i) Uji rasio itu syarat cukup, bukan syarat perlu:} ketika
$\abs{a_{n+1}/a_n}$ tak berlimit
(\cref{ex:b2:powerseries:sinn}, \cref{exo:b2:powerseries:1}),
kembalilah ke definisinya: $R = \sup\{r : (a_nr^n)$
terbatas$\}$. \emph{(ii) Tak ada yang menyeberangi perbatasannya secara
cuma-cuma:} penurunan dan pengintegralan suku demi suku adalah teorema
\emph{di dalam} cakram \omterm{def:b2:metric:topology}{terbukanya}; sedangkan di $\abs x = R$ tiap
deretnya wajib diperiksa ulang (dan itulah seluruh bahasan soal akhir
pekan). \emph{(iii) Jari-jari sebuah jumlah:} $\min(R_a, R_b)$
hanyalah batas bawah --- sebab pencoretan dapat membesarkannya ($a_n =
1, b_n = -1$: jumlahnya $0$ secara identik, jari-jarinya $\infty$).
\emph{(iv) $C^\infty$ bukan \omterm{def:b2:powerseries:analytic}{analitik}:} deret Taylor yang konvergen
boleh jadi konvergen ke fungsi yang \emph{salah}
(\cref{exo:b2:powerseries:7}); jadi sebelum menulis $f(x) = \sum
\frac{f^{(n)}(0)}{n!}x^n$, buktikanlah dulu --- lewat sebuah persamaan
diferensial (\cref{met:b2:powerseries:ode}), taksiran sisanya, atau
sebuah rumus integral.
\end{remark}

\begin{remark}[Di mana ini dipakai]
Deret pangkat menjadi kuda beban tiga bab berikutnya: bab
persamaan diferensial menyelesaikan persamaan linear dengan menyuntikkan
$\sum a_nx^n$ (yakni kotak metode di bawah, yang diindustrikan); bab
\omterm{ex:b2:powerseries:fibonacci}{fungsi pembangkit} mengubah kesamaan tentang
peluang menjadi kesamaan tentang deret pangkat lalu kembali lagi; dan
jilid Tahun ke-3 meresmikan peubah kompleksnya, tempat
keanalitikan menjadi setara dengan keterdiferensialan kompleks dan
``pemusatan ulang yang diterima'' di atas mendapat buktinya yang jujur.
Adapun soal akhir pekan menjelajahi satu tempat yang didiamkan teorema
bab ini: yakni perbatasan $\abs x = R$ itu sendiri.
\end{remark}

\begin{method}[Menguraikan lewat sebuah persamaan diferensial]\label{met:b2:powerseries:ode}
Untuk menguraikan sebuah fungsi $f$ menjadi deret pangkat: carilah
persamaan diferensial linear berkoefisien polinomial yang dipenuhi $f$;
suntikkan $\sum a_nx^n$; samakan koefisiennya untuk memperoleh rekurensi
bagi $(a_n)$; selesaikan, lalu periksa jari-jari dan syarat awalnya.
Contohnya --- deret binomial: $f(x) = (1+x)^\alpha$ memenuhi $(1+x)f' =
\alpha f$ dan $f(0) = 1$; menyuntikkannya memberi $(n+1)a_{n+1} =
(\alpha - n)a_n$, jadi $a_n = \binom{\alpha}{n}$, berjari-jari $1$
(lewat uji rasio), dan jumlahnya, yang memenuhi persamaan diferensial
yang sama dengan nilai awal yang sama, sama dengan $(1 + x)^\alpha$
berkat teorema ketunggalan bagi persamaan diferensial linear (jilid
Tahun ke-1).
\end{method}

\begin{example}[Metodenya pada sebuah persamaan berpaksa]\label{ex:b2:powerseries:odeforced}
Selesaikan $y' = y + x$, $y(0) = 0$, lewat deret pangkat. Dengan
menyuntikkan $y = \sum a_nx^n$:
\[
\sum_{n\geq0}(n+1)a_{n+1}x^n
= \sum_{n\geq0}a_nx^n + x ,
\]
lalu menyamakan koefisiennya: $a_1 = a_0 = 0$, $2a_2 = a_1 + 1
= 1$, dan $(n+1)a_{n+1} = a_n$ untuk $n \geq 2$. Jadi $a_2 =
\frac{1}{2!}$ dan, secara induktif, $a_n = \frac{1}{n!}$ untuk setiap
$n \geq 2$: jari-jarinya $\infty$, dan
\[
y(x) = \sum_{n\geq2}\frac{x^n}{n!} = \eu^x - 1 - x .
\]
Periksa: $y' = \eu^x - 1 = y + x$ dan $y(0) = 0$. Pelajaran
penutupnya: rekurensinya \emph{adalah} persamaannya, koefisien demi
koefisien; sedangkan suku paksanya hanya mengusik berhingga banyak
koefisien awal, dan sesudah itu pola homogennya mengambil
alih --- yakni bayangan diskret bagi ``penyelesaian khusus ditambah
penyelesaian homogen''.
\end{example}

\section{Fungsi pembangkit}

\begin{example}[Fibonacci]\label{ex:b2:powerseries:fibonacci}
Misalkan $F(x) = \sum_{n\geq0} F_n x^n$ (dengan bilangan Fibonacci $F_0
= 0$, $F_1 = 1$). Rekurensi $F_{n+2} = F_{n+1} + F_n$ diterjemahkan,
setelah dikalikan $x^{n+2}$ lalu dijumlahkan, menjadi
\[
F(x) - x = x\,F(x) + x^2 F(x)
\quad\Longrightarrow\quad
F(x) = \frac{x}{1 - x - x^2} ,
\]
yang sahih di tempat deretnya konvergen. Jari-jarinya
$\frac{1}{\varphi}$: sebab dari $F_n \sim
\frac{\varphi^n}{\sqrt5}$ (yakni Binet pada contoh berikutnya --- atau
lewat induksi kasar $F_n \leq 2^n$ ditambah rekurensinya), uji rasio
memberi
\[
\frac{F_{n+1}\abs x^{n+1}}{F_n\abs x^n}
\longrightarrow \varphi\abs x ,
\qquad\text{konvergen bila dan hanya bila } \abs x < \frac1\varphi
\approx 0.618 .
\] Pecahan parsial pada
$\frac{x}{1 - x - x^2}$ beserta deret geometrinya menurunkan ulang rumus
Binet --- jadi \omterm{ex:b2:powerseries:fibonacci}{fungsi pembangkit} mengindustrikan rekurensi
linear.\index{fungsi pembangkit}
\end{example}

\begin{example}[Rumus Binet, dikerjakan]\label{ex:b2:powerseries:binet}
Misalkan $\varphi = \frac{1+\sqrt5}{2}$ dan $\psi =
\frac{1-\sqrt5}{2}$, yakni akar $X^2 = X + 1$; karena $\varphi
+ \psi = 1$ dan $\varphi\psi = -1$, berlaku
\[
1 - x - x^2 = (1 - \varphi x)(1 - \psi x) .
\]
Pecahan parsialnya: dengan mencari $\frac{x}{(1-\varphi x)(1-\psi x)} =
\frac{A}{1 - \varphi x} + \frac{B}{1 - \psi x}$, suku tetapnya
memberi $A + B = 0$ dan koefisien $x$-nya $-A\psi - B\varphi = 1$,
sehingga $A(\varphi - \psi) = 1$: jadi $A = \frac{1}{\sqrt5} = -B$.
Dua deret geometri kemudian,
\[
F(x) = \frac{1}{\sqrt5}\sum_{n\geq0}
\bigl(\varphi^n - \psi^n\bigr)x^n
\quad\Longrightarrow\quad
F_n = \frac{\varphi^n - \psi^n}{\sqrt5}
\]
berkat ketunggalan koefisiennya
(\cref{thm:b2:powerseries:calculus}). Karena $\abs\psi < 1$, suku
$\frac{\psi^n}{\sqrt5}$ bernilai mutlak $< \frac12$: jadi
$F_n$ adalah \emph{bilangan bulat terdekat} dari
$\frac{\varphi^n}{\sqrt5}$. Pelajaran penutupnya: jari-jari
$\frac1\varphi$ milik $F$ adalah kebalikan akar dominannya
--- jadi pertumbuhan koefisien dan \omterm{def:b2:powerseries:radius}{jari-jari kekonvergenan} adalah
informasi yang sama, dibaca dari arah yang berlawanan.
\end{example}

\begin{example}[Bilangan Catalan]\label{ex:b2:powerseries:catalan}
\omterm{ex:b2:powerseries:catalan}{Bilangan Catalan} $C_n$ (yakni banyaknya triangulasi, pengurungan,
lintasan Dyck, \dots) memenuhi $C_0 = 1$ dan $C_{n+1} = \sum_{k=0}^n
C_kC_{n-k}$. \omterm{ex:b2:powerseries:fibonacci}{Fungsi pembangkitnya} $C(x) = \sum C_nx^n$ lalu
memenuhi (lewat hasil kali Cauchy!)
\[
C(x) = 1 + x\,C(x)^2
\quad\Longrightarrow\quad
C(x) = \frac{1 - \sqrt{1 - 4x}}{2x} ,
\]
dengan memilih akar yang $C(0) = 1$: sebab menyelesaikan persamaan
kuadrat $xC^2 - C + 1 = 0$ memberi dua calon $\frac{1 \pm
\sqrt{1-4x}}{2x}$, dan ketika $x \to 0$ akar ``$+$'' meledak
seperti $\frac1x$ sedangkan akar ``$-$'' menuju $1$ (uraikan
$\sqrt{1-4x} = 1 - 2x + O(x^2)$) --- jadi hanya tanda minusnya yang
dapat mengusung deret pangkat dengan $C_0 = 1$. Menguraikan $\sqrt{1 -
4x}$ lewat deret binomial memberi bentuk tertutupnya
\[
C_n = \frac{1}{n+1}\binom{2n}{n} ,
\]
yang dikerjakan pada \cref{exo:b2:powerseries:8}.\index{bilangan Catalan}
\end{example}

\begin{remark}[Deret formal lawan deret konvergen]
Setiap perhitungan \omterm{ex:b2:powerseries:fibonacci}{fungsi pembangkit} di atas berakhir dengan memanggil
ketunggalan koefisiennya, dan teorema itu hidup
\emph{di dalam} cakram berjari-jari positif: jadi sebelum ``membaca''
$F_n$ atau $C_n$, kita wajib mengetahui $R > 0$. Batas apriori yang
kasar sudah cukup --- $F_n \leq 2^n$ (lewat induksi seketika) memberi $R
\geq \frac12$ bagi Fibonacci; $C_n \leq 4^n$ (sebab tiap \omterm{ex:b2:powerseries:catalan}{bilangan
Catalan} mencacah himpunan bagian lintasan) memberi $R \geq \frac14$.
Awas ujung skalanya yang merosot: $\sum n!\,x^n$ berjari-jari $0$, dan
memperlakukannya sebagai fungsi tak bermakna --- sebab kesamaan
yang melibatkan deret semacam itu menjadi milik kalkulus koefisien yang
\emph{formal}, yakni permainan aljabar murni dengan aturannya sendiri
(yang berbeda). Pada tingkat ini: amankanlah dulu jari-jari yang positif,
lalu hitunglah sebebas-bebasnya di dalamnya.
\end{remark}

\begin{remark}[Pandangan ke depan di dalam jilid ini]
Deret pangkat adalah satu dari dua mesin uraian besar buku ini;
satunya adalah deret Fourier pada bab harmonik,
dan membandingkan keduanya sangat mendidik. Deret pangkat bersifat
kaku: koefisiennya terpaksa ($a_n = f^{(n)}(0)/n!$),
kekonvergenannya kejam (normal di dalam, sia-sia di luar),
dan jumlahnya \omterm{def:b2:powerseries:analytic}{analitik} --- yakni kaku tanpa batas
(\cref{def:b2:powerseries:analytic}). Sedangkan deret Fourier bersifat
lentur: ia menyajikan sinyal yang sekadar mulus sepotong-sepotong, dengan
harga berupa pertanyaan kekonvergenan yang peka di perbatasan
kemulusannya. Kedua teori itu bertemu pada soal akhir pekan bab ini:
yakni penjumlahan Cesàro dan Abel, yang dikembangkan di sini untuk
lingkaran batasnya, lalu kembali pada bab Fourier sebagai kernel Fejér
dan kernel Poisson. Sementara itu bab persamaan diferensial
mengonsumsi deret pangkat secara langsung ($\eu^{tA}$, penyelesaian
berupa deret), dan bab \omterm{ex:b2:powerseries:fibonacci}{fungsi pembangkit} mengubah
siasat \cref{ex:b2:powerseries:fibonacci} menjadi kalkulus
yang sistematis bagi peluang.
\end{remark}

\section{Latihan}

\begin{exercise}[$\star$]\label{exo:b2:powerseries:1}
\omterm{def:b2:powerseries:radius}{Jari-jari kekonvergenannya}:
$\sum \dfrac{n^2}{2^n}z^n$; $\;\sum \dfrac{z^n}{\binom{2n}{n}}$;
$\;\sum z^{n!}$; $\;\sum \bigl(2 + (-1)^n\bigr)^n z^n$.
\end{exercise}

\begin{exercise}[$\star$]\label{exo:b2:powerseries:2}
Tunjukkan bahwa $\sum z^n$, $\sum \frac{z^n}{n}$, $\sum \frac{z^n}{n^2}$
semuanya berjari-jari $1$ tetapi berkelakuan berbeda di $z = 1$ dan $z =
-1$: berturut-turut divergen/divergen, divergen/konvergen, dan
konvergen/konvergen.
\end{exercise}

\begin{exercise}[$\star$]\label{exo:b2:powerseries:3}
Hitunglah jumlahnya, untuk $\abs x < 1$:
\[
\sum_{n\geq0} n x^n,
\qquad
\sum_{n\geq0} n^2 x^n,
\qquad
\sum_{n\geq0} \frac{x^{2n+1}}{2n+1} .
\]
\end{exercise}

\begin{exercise}[$\star\star$]\label{exo:b2:powerseries:4}
Uraikan menjadi deret pangkat di $0$, beserta jari-jarinya:
$\dfrac{1}{(1-x)(2-x)}$ (lewat pecahan parsial); dan $\;\ln(1 + x +
x^2)$ \emph{(tulis $1 + x + x^2 = \frac{1 - x^3}{1 - x}$)}.
\end{exercise}

\begin{exercise}[$\star\star$]\label{exo:b2:powerseries:5}
Buktikan bahwa $f(x) = \sum_{n\geq1} H_n x^n = -\dfrac{\ln(1 -
x)}{1 - x}$ untuk $\abs x < 1$, dengan $H_n$ menyatakan bilangan harmonik
\emph{(lewat hasil kali Cauchy $\sum x^n$ dan $\sum \frac{x^n}{n}$)}.
\end{exercise}

\begin{exercise}[$\star\star$]\label{exo:b2:powerseries:6}
Selesaikan lewat \omterm{ex:b2:powerseries:fibonacci}{fungsi pembangkit} rekurensi $u_0 = 1$, $u_{n+1} =
2u_n + n$: hitunglah $U(x) = \sum u_nx^n$ dalam bentuk tertutup, uraikan,
lalu bacalah $u_n = 2^{n+1} - n - 1$.
\end{exercise}

\begin{exercise}[$\star\star$]\label{exo:b2:powerseries:7}
Misalkan $f(x) = \eu^{-1/x^2}$ untuk $x \neq 0$ dan $f(0) = 0$.
Buktikan bahwa $f$ berkelas $C^\infty$ pada $\R$ dengan $f^{(n)}(0) = 0$
untuk setiap $n$ \emph{(tunjukkan secara induktif bahwa $f^{(n)}(x) =
P_n\bigl(\frac1x\bigr) \eu^{-1/x^2}$ untuk suatu polinomial $P_n$, lalu
pakai perbandingan pertumbuhannya)}. Simpulkan bahwa $f$ tidak \omterm{def:b2:powerseries:analytic}{analitik}
di $0$: sebab deret Taylornya di $0$ konvergen --- ke fungsi yang salah.
\end{exercise}

\begin{exercise}[$\star\star\star$]\label{exo:b2:powerseries:8}
Lengkapilah \cref{ex:b2:powerseries:catalan}: uraikan $\sqrt{1 - 4x}$
lewat deret binomial, dengan menunjukkan
\[
\binom{1/2}{n+1}(-4)^{n+1} = -\frac{2}{n+1}\binom{2n}{n},
\]
lalu turunkan $C_n = \frac{1}{n+1}\binom{2n}{n}$; tentukan pula
\omterm{def:b2:powerseries:radius}{jari-jari kekonvergenan} $C(x)$ dan asimtotik $C_n$ lewat
Stirling.
\end{exercise}

\begin{exercise}[$\star\star\star$]\label{exo:b2:powerseries:9}
(Teorema limit radial Abel, kasus khusus) Andaikan $\sum a_n$
konvergen. Buktikan bahwa $\lim_{x \to 1^-} \sum_{n} a_n x^n = \sum_n
a_n$. \emph{(Lewat \omterm{thm:b2:series:abel}{penjumlahan Abel}: dengan $A_n$ jumlah parsialnya dan
$A = \lim A_n$, tulislah $\sum a_nx^n = (1 - x)\sum A_n x^n$; lalu $\sum
a_nx^n - A = (1-x)\sum (A_n - A)x^n$, lalu pecah jumlahnya di $N$ yang
besar.)} Penerapannya: $\sum \frac{(-1)^{n-1}}{n} = \ln 2$ dan $\sum
\frac{(-1)^n}{2n+1} = \frac\pi4$, dibuktikan ulang dari deret pangkatnya.
\end{exercise}

\begin{exercise}[$\star$]\label{exo:b2:powerseries:10}
Tunjukkan bahwa $\displaystyle\sum_{n\geq1}\frac{x^n}{n(n+1)} = 1 +
\frac{1-x}{x}\,\ln(1-x)$ untuk $0 < \abs x < 1$, tentukan
jari-jarinya, lalu periksalah bahwa kekonvergenannya normal pada
$\intcc{-1}{1}$; periksa pula bahwa nilai di $x = 1$ yang diramalkan
kekontinuannya sepakat dengan jumlah teleskop $\sum
\frac{1}{n(n+1)} = 1$.
\end{exercise}

\begin{exercise}[$\star\star$]\label{exo:b2:powerseries:11}
(Permutasi tanpa titik tetap) Misalkan $D_n$ banyaknya permutasi atas
$n$ objek yang tak punya titik tetap ($D_0 = 1$). Menyortir
permutasi $\{1, \dots, n\}$ menurut himpunan titik tetapnya memberi
$n! = \sum_{k=0}^{n}\binom nk D_{n-k}$. Kalikan dengan
$\frac{x^n}{n!}$, jumlahkan, lalu kenali sebuah hasil kali Cauchy untuk
memperoleh \omterm{ex:b2:powerseries:fibonacci}{fungsi pembangkit} eksponensialnya
\[
\sum_{n\geq0} D_n\,\frac{x^n}{n!} = \frac{\eu^{-x}}{1-x}
\qquad (\abs x < 1),
\]
lalu bacalah bentuk tertutupnya $\dfrac{D_n}{n!} =
\sum_{k=0}^{n}\dfrac{(-1)^k}{k!}$ beserta limitnya $\dfrac{D_n}{n!}
\to \eu^{-1}$.
\end{exercise}

\begin{exercise}[$\star\star\star$]\label{exo:b2:powerseries:12}
Buktikan, dengan deret binomial pada
\cref{met:b2:powerseries:ode}, bahwa
\[
\frac{1}{\sqrt{1 - 4x}} = \sum_{n\geq0}\binom{2n}{n}x^n
\qquad \Bigl(\abs x < \frac14\Bigr),
\]
lalu turunkan, dengan mengkuadratkannya (yakni hasil kali Cauchy
terhadap $\frac{1}{1-4x} = \sum 4^nx^n$), kesamaan konvolusinya
\[
\sum_{k=0}^{n}\binom{2k}{k}\binom{2n-2k}{n-k} = 4^n .
\]
\end{exercise}

\section{Soal: Abel, Tauber, dan perbatasan kekonvergenan}

\begin{problem}\label{pb:b2:powerseries:1}
Di dalam cakram kekonvergenannya segalanya mudah; seluruh
dramanya justru terjadi \emph{di} perbatasan. Soal ini membangun teori
perbatasan itu dalam peubah real: teorema Abel dalam bentuk
\omterm{def:b2:funcseq:def}{seragamnya}, konversnya di bawah syarat Tauber, hierarki metode
penjumlahan Cesàro--Abel (beserta teorema Frobenius),
pengintegralan suku demi suku sampai ke perbatasannya dengan
konstanta klasik sebagai panennya, dan akhirnya kekakuan \omterm{def:b2:powerseries:analytic}{fungsi
analitik} --- yakni teorema
identitas.\index{teorema Abel}\index{teorema Tauber}
Di sepanjang soal ini, $(a_n)$ barisan real, $f(x) = \sum_{n\geq0}
a_nx^n$, dan $A_n = a_0 + \dots + a_n$.

\medskip
\noindent\textbf{Bagian I --- Teorema Abel, secara \omterm{def:b2:funcseq:def}{seragam}.} Pada
bagian ini andaikan $\sum a_n$ konvergen, lalu tetapkan $r_n =
\sum_{k\geq n} a_k$ (sehingga $r_n \to 0$ dan $a_n = r_n - r_{n+1}$).
\begin{enumerate}
  \item Buktikan, lewat penjumlahan parsial, bahwa untuk setiap $0 \leq
        x \leq 1$ dan $N \leq M$:
        \[
        \Bigl|\sum_{n=N}^{M} a_n x^n\Bigr|
        \leq 2\sup_{n \geq N}\,\abs{r_n} .
        \]
  \item Turunkan bahwa $\sum a_nx^n$ konvergen \emph{\omterm{def:b2:funcseq:def}{seragam}} pada
        $\intcc{0}{1}$, bahwa jumlahnya \omterm{def:b2:metric:continuity}{kontinu} di sana, lalu
        peroleh kembali limit radial pada
        \cref{exo:b2:powerseries:9}: $f(x) \to \sum a_n$ ketika $x
        \to 1^-$.
  \item (Teorema Abel untuk hasil kali Cauchy) Misalkan $\sum a_n = A$,
        $\sum b_n = B$ dan andaikan hasil kali Cauchy $\sum
        c_n$, dengan $c_n = \sum_{k} a_kb_{n-k}$, \emph{konvergen}
        dengan jumlah $C$. Buktikan $C = AB$ \emph{(sebab di dalam
        cakramnya kesamaan hasil kalinya berlaku menurut
        \cref{prop:b2:powerseries:operations}; lalu lewatkan $x \to
        1^-$)}.
  \item Tunjukkan bahwa hipotesisnya penting: untuk $a_n = b_n =
        \frac{(-1)^n}{\sqrt{n+1}}$, kedua deretnya konvergen, namun
        $\abs{c_n} \geq \frac{2(n+1)}{n+2} \geq 1$ \emph{(batasilah
        tiap faktor $\sqrt{(k+1)(n-k+1)}$ lewat AM--GM)}: jadi
        hasil kali Cauchy dua deret konvergen boleh saja divergen.
  \item (Sebuah panen \cref{exo:b2:powerseries:5}) Tunjukkan bahwa
        $\bigl(\ln(1-x)\bigr)^2 = 2\sum_{n\geq1}
        \frac{H_n}{n+1}x^{n+1}$ pada $\intoo{-1}{1}$, periksalah bahwa
        $\bigl(\frac{H_n}{n+1}\bigr)_{n\geq1}$ turun ke
        $0$, lalu rampungkan dengan Abel:
        \[
        \sum_{n\geq1} (-1)^{n+1}\,\frac{H_n}{n+1}
        = \frac{(\ln 2)^2}{2} .
        \]
\end{enumerate}

\medskip
\noindent\textbf{Bagian II --- Konvers Tauber.} Sebutlah $\sum
a_n$ \emph{terjumlahkan-Abel} ke $L$ apabila $f(x) \to L$ ketika $x \to
1^-$.
\begin{enumerate}[resume]
  \item Tunjukkan bahwa $\sum (-1)^n$ terjumlahkan-Abel ke $\frac12$
        padahal divergen: jadi teorema Abel tak punya
        konvers tanpa syarat.
  \item (Lema Cesàro) Jika $u_n \to 0$ maka $\frac{u_1 + \dots
        + u_N}{N} \to 0$ \emph{(pecahlah jumlahnya di sebuah $m$ yang
        tetap)}.
  \item Kini andaikan $n\,a_n \to 0$ dan $f(x) \to L$. Dengan $x_N =
        1 - \frac1N$, buktikan kedua taksiran
        \[
        \Bigl|\sum_{n=0}^{N} a_n\bigl(1 - x_N^n\bigr)\Bigr|
        \leq \frac{1}{N}\sum_{n=1}^{N} n\,\abs{a_n},
        \qquad
        \Bigl|\sum_{n>N} a_n x_N^n\Bigr|
        \leq \sup_{n>N}\bigl(n\abs{a_n}\bigr)
        \]
        \emph{(untuk yang pertama, $1 - x^n \leq n(1-x)$; untuk yang
        kedua, $\abs{a_n} \leq \frac{1}{N}\sup_{m>N} m\abs{a_m}$ dan
        $\sum x_N^n \leq N$)}.
  \item Simpulkan \emph{teorema Tauber}: jika $n\,a_n \to 0$ dan
        $\sum a_n$ terjumlahkan-Abel ke $L$, maka $\sum a_n$
        konvergen ke $L$.
  \item (Tauberian yang mudah untuk koefisien positif) Jika $a_n
        \geq 0$ dan $f$ terbatas pada $\intco{0}{1}$, tunjukkan bahwa
        $\sum a_n$ konvergen dan $\sum a_n = \lim_{x\to1^-}
        f(x)$ \emph{(batasi $\sum_{n\leq N}a_nx^n \leq f(x)$ lalu
        lewatkan $x \to 1^-$, kemudian pakai Abel)}.
\end{enumerate}

\medskip
\noindent\textbf{Bagian III --- Rata-rata Cesàro dan teorema
Frobenius.} Sebutlah $\sum a_n$ \emph{terjumlahkan-Cesàro} ke $L$
apabila $\sigma_N = \frac{A_0 + \dots + A_{N-1}}{N} \to L$.
\begin{enumerate}[resume]
  \item Tunjukkan bahwa deret yang konvergen bersifat
        terjumlahkan-Cesàro ke jumlahnya \emph{(pertanyaan 7 yang
        diterapkan pada $A_n - L$)}.
  \item Hitunglah nilai Cesàro bagi $\sum(-1)^n$ lalu periksa bahwa ia
        sepakat dengan nilai Abelnya $\frac12$ pada pertanyaan 6.
  \item Dengan $S_n = A_0 + \dots + A_n = (n+1)\,\sigma_{n+1}$,
        buktikan kedua kesamaan berikut, untuk $0 \leq x < 1$:
        \[
        f(x) = (1-x)^2\sum_{n\geq0}(n+1)\,\sigma_{n+1}x^n,
        \qquad
        (1-x)^2\sum_{n\geq0}(n+1)x^n = 1 .
        \]
  \item (Frobenius) Turunkan: jika $\sigma_N \to L$ maka $f(x) \to
        L$ ketika $x \to 1^-$ --- jadi terjumlahkan-Cesàro
        mengakibatkan terjumlahkan-Abel, dengan nilai yang sama
        \emph{(kurangkan kedua kesamaannya lalu pecah jumlahnya di $N$
        yang besar, seperti pada \cref{exo:b2:powerseries:9})}.
  \item Tunjukkan bahwa hierarki
        \[
        \text{konvergen} \;\Longrightarrow\;
        \text{terjumlahkan-Cesàro} \;\Longrightarrow\;
        \text{terjumlahkan-Abel}
        \]
        bersifat tegas pada kedua panahnya: pertanyaan 6 untuk yang
        pertama; sedangkan untuk yang kedua, tunjukkan bahwa
        $\sum(-1)^n(n+1)$ terjumlahkan-Abel ke $\frac14$ (hitunglah
        $f$) tetapi tidak terjumlahkan-Cesàro (hitunglah $\sigma_N$
        secara terpisah untuk $N$ genap dan untuk yang ganjil).
\end{enumerate}

\medskip
\noindent\textbf{Bagian IV --- Mengintegralkan sampai ke perbatasan.}
\begin{enumerate}[resume]
  \item Andaikan $\sum a_nx^n$ konvergen pada $\intco{0}{1}$ dan
        $\sum \frac{a_n}{n+1}$ konvergen. Buktikan bahwa
        \omterm{def:b2:integration:improper}{integral tak wajar} $\int_0^1 f$ ada dan
        \[
        \int_0^1 \Bigl(\sum_{n\geq0} a_nx^n\Bigr)\dd x
        = \sum_{n\geq0}\frac{a_n}{n+1}
        \]
        \emph{(sebab antiturunannya $F(x) =
        \sum\frac{a_n}{n+1}x^{n+1}$ \omterm{def:b2:metric:continuity}{kontinu} di $1$ menurut Bagian I)}.
  \item Misalkan $\eta = \sum_{n\geq1}\frac{(-1)^{n-1}}{n^2}$.
        Tunjukkan $\int_0^1 \frac{\ln(1+x)}{x}\dd x = \eta$ dan,
        dengan memisahkan indeks genap dan ganjil pada $\sum
        \frac1{n^2}$ yang konvergen \omterm{def:b2:series:def}{mutlak}, bahwa $\eta =
        \frac12\sum_{n\geq1}\frac{1}{n^2}$. (Soal akhir pekan bab
        Fourier menilai $\sum\frac1{n^2} = \frac{\pi^2}{6}$.)
  \item Buktikan
        \[
        \sum_{n\geq0}\frac{(-1)^n}{3n+1}
        = \int_0^1\frac{\dd x}{1+x^3}
        = \frac13\Bigl(\ln 2 + \frac{\pi}{\sqrt3}\Bigr)
        \]
        \emph{(deretnya konvergen menurut Leibniz; integralkan
        deret geometri $\sum(-1)^nx^{3n}$ dengan pertanyaan 16;
        lalu pecahan parsial: $\frac{1}{1+x^3} =
        \frac{1/3}{1+x} + \frac{(2-x)/3}{x^2-x+1}$)}.
  \item Dari deret binomial bagi $(1-t)^{-1/2}$
        (\cref{exo:b2:powerseries:12}) turunkan
        \[
        \arcsin x = \sum_{n\geq0}
        \frac{\binom{2n}{n}}{4^n(2n+1)}\,x^{2n+1}
        \quad(\abs x < 1),
        \qquad\text{lalu}\qquad
        \sum_{n\geq0}\frac{\binom{2n}{n}}{4^n(2n+1)}
        = \frac\pi2 ,
        \]
        dengan membenarkan nilai perbatasannya lewat kekonvergenan
        \emph{normal} pada $\intcc{-1}{1}$ (pakailah
        $\binom{2n}n4^{-n} \sim \frac{1}{\sqrt{\pi n}}$,
        \cref{ex:b2:comparison:centralbinomial}) --- di sini bahkan
        Abel tidak diperlukan.
  \item (Catalan di perbatasannya) Tunjukkan bahwa $\sum C_n 4^{-n} =
        2$: jadi deret Catalan pada
        \cref{ex:b2:powerseries:catalan} konvergen \emph{di}
        jari-jarinya $\frac14$ (lewat asimtotik
        \cref{exo:b2:powerseries:8}), jumlahnya \omterm{def:b2:metric:continuity}{kontinu} pada
        $\intcc{0}{\frac14}$, dan bentuk tertutupnya berlimit $2$
        di sana.
\end{enumerate}

\medskip
\noindent\textbf{Bagian V --- Kekakuan: teorema identitas.}
\begin{enumerate}[resume]
  \item (Nol yang terpencil) Misalkan $f = \sum a_nx^n$ berjari-jari $R
        > 0$ dan tidak semua $a_n = 0$; misalkan $m$ indeks terkecil
        dengan $a_m \neq 0$. Tunjukkan $f(x) = x^m g(x)$ dengan $g$
        sebuah deret pangkat berjari-jari $R$, $g(0) = a_m \neq 0$,
        lalu turunkan bahwa $f$ tak punya nol pada suatu
        persekitaran $0$ yang tertusuk.
  \item (Teorema identitas) Misalkan $f, h$ jumlah deret pangkat
        di dekat $0$ dan $(x_k)$ barisan titik \emph{tak nol}
        dengan $x_k \to 0$ dan $f(x_k) = h(x_k)$. Buktikan
        bahwa $f$ dan $h$ punya koefisien yang sama, sehingga
        keduanya berimpit di dekat $0$.
  \item Carilah \emph{semua} fungsi $f$ yang \omterm{def:b2:powerseries:analytic}{analitik} di dekat $0$ dengan
        \[
        f\Bigl(\frac1k\Bigr) = \frac{k^2}{k^2+1}
        \qquad\text{untuk setiap bilangan bulat besar } k .
        \]
  \item Tunjukkan bahwa \omterm{def:b2:powerseries:analytic}{fungsi analitik} pada interval \omterm{def:b2:metric:topology}{terbuka} $I$
        yang nol pada sebuah subinterval pastilah nol secara identik pada
        $I$ \emph{(sebab himpunan titik yang di sekitarnya $f$ nol
        secara identik bersifat \omterm{def:b2:metric:topology}{terbuka} dan, lewat teorema identitas
        yang diterapkan di titik akumulasinya, tertutup di $I$)}.
        Simpulkan bahwa tak ada \omterm{def:b2:powerseries:analytic}{fungsi analitik} tak nol pada $\R$ yang
        berpenyangga \omterm{def:b2:metric:compact}{kompak} --- padahal fungsi bonggol $C^\infty$
        memang ada (\cref{exo:b2:powerseries:7} menyediakan
        batu bangunannya): jadi keanalitikan itu kaku, kemulusan itu
        lembek.
  \item Rangkuman. Satu kalimat untuk masing-masing: (i) apa yang
        ditambahkan teorema Abel pada paket \omterm{def:b2:funcseq:series}{kekonvergenan normal} pada
        \cref{lem:b2:powerseries:abel}; (ii) hipotesis persis yang
        membuat konversnya berlaku (Tauber) dan
        anak tangga antaranya (Frobenius); (iii) satu konstanta
        perbatasan dari Bagian IV yang kini dapat kaupaparkan kepada
        seorang teman dalam dua baris; (iv) di mana rata-rata Cesàro
        akan muncul kembali pada buku ini, untuk deret berjenis yang
        sangat berbeda.

\end{enumerate}
\end{problem}
